\chapter{Introduccion}
\noindent
Antes de iniciar con el contenido teórico de la asignatura, debemos recalcar varios puntos a tener en cuenta. \\

\noindent
El primero y más importante, aunque los apuntes estén en español, la asignatura está totalmente impartida en italiano, y el examen escrito está en italiano, por lo que es conveniente que se tenga en cuenta la traducción de ciertos conceptos y términos importantes, bien sean por ejemplo, \textit{spazio vettoriale normato} (espacio vectorial normado), \textit{spazio di Banach} (espacio de Banach), \textit{spazio di Hilbert} (espacio de Hilbert), \textit{operatore lineare} (operador lineal), \textit{operatore limitato} (operador acotado), \textit{operatore compatto} (operador compacto), \textit{operatore autoaggiunto} (operador autoadjunto), \textit{operatore unitario} (operador unitario), \textit{operatore normale} (operador normal), \textit{spettro} (espectro), etc. \\

\noindent
La profesora que impartió la asignatura fue \tbf{Filomena Pacella}, con ella aunque el examen escrito esté en italiano, puedes resolverlo comentando lo que necesites en inglés, y además puedes hacer el examen oral en inglés. \\ \\

\noindent
El segundo punto a recalcar es que la nota final no se basa en la media aritmética entre el examen oral y escrito, sino que primero se hace la prueba escrita (que necesita ser aprobada para hacer el examen oral), y después se hace el examen oral, donde ella pone la nota final basándose gran parte en el examen oral, aunque como es obvio el escrito también tiene peso si se obtuvo con buena calificación. \\ \\

\noindent
Por último, como consejo para el examen oral, es conveniente tener control sobre la mayoría de definiciones, teoremas y resultados, pero lo que más se va a preguntar es la demostración de los teoremas más importantes de la asignatura (aquellos con nombre propio), como por ejemplo el teorema de Hahn-Banach, el teorema de Banach-Steinhaus, el teorema de Open Mapping, el teorema de Closed Graph, el teorema de Riesz Representation, etc. \\ \\